\appendix
\chapter{问答与讨论}

\section{Tang \& Rong 系列论文核心概念}

本节记录学习 Tang \& Rong 信息几何系列论文过程中的问答与思考，涵盖流形推断、扩散模型、贝叶斯稳健推断、MCMC 等核心主题。

\subsection{什么是黎曼流形上的分布推断}\label{sec:qa-rimfd}

\textbf{问}：什么是黎曼流形上的分布推断（Distribution Submanifold）？

\textbf{答}：分布子流形是指概率分布空间的一个低维嵌入子流形。在指数族中，分布可以用充分统计量的均值参数化，这些均值向量构成一个线性子空间与概率单纯形的交集——本质上就是一个统计流形。

Tang (2022) 的核心贡献是证明：当分布族是一个 $d$ 维光滑子流形时，最优的分布估计率由曲率决定，且满足：
$$D(P_{\hat{\theta}_n} \| P_{\theta^*}) \asymp \frac{d}{n} + \kappa^2 \cdot \text{Curvature}(M)$$

其中 $\kappa$ 是子流形的曲率尺度。这将几何与统计推断的核心联系了起来。

\subsection{为什么曲率会影响估计效率}\label{sec:qa-curvature-efficiency}

\textbf{问}：为什么子流形的曲率会影响统计推断的效率？

\textbf{答}：曲率反映了流形的"弯曲程度"，它从两个渠道影响估计：

\begin{enumerate}
    \item \textbf{近似误差}：用参数模型 $M$ 逼近真实分布 $P^*$ 时，真实分布与最近模型分布的距离与曲率成正比。高曲率意味着更大的近似偏差。
    \item \textbf{方差放大}：在曲率大的区域，参数空间的小扰动会导致分布空间的较大变化，使估计方差增大。
\end{enumerate}

几何直觉：平坦空间中的"两点之间直线最短"在弯曲空间中不再成立——测地线需要绕弯，这增加了估计的不确定性。

\subsection{扩散模型如何学习黎曼度量}\label{sec:qa-diffusion-riemannian-metric}

\textbf{问}：扩散模型（Score-based models）如何学习数据流形上的黎曼度量？

\textbf{答}：Tang (2024) 的 Diffusion Manifold Adaptivity 提出了一个优雅的框架：

\begin{enumerate}
    \item \textbf{Score 估计}：训练神经网络 $s_\theta(x, t) \approx \nabla_x \log p_t(x)$ 估计 Fisher score
    \item \textbf{Riemannian metric 学习}：通过分析 score 的协方差结构，估计黎曼度量张量 $g_{ij}(x)$
    \item \textbf{几何一致性}：确保学习到的度量与数据的几何结构一致
\end{enumerate}

关键洞察：扩散模型的去噪过程天然地与流形的几何结构相关联——在低噪声水平下，score 指向流形的局部切空间方向。

\subsection{什么是流形上的稳健贝叶斯推断}\label{sec:qa-robust-bayesian-manifold}

\textbf{问}：什么是黎曼流形上的稳健贝叶斯推断？与欧氏空间有何不同？

\textbf{答}：Robust Bayesian Inference on Riemannian Manifolds (Tang 2023) 研究的核心问题是：当数据中存在污染或异常值时，如何在流形上进行稳健的参数估计。

关键差异在于：

\begin{itemize}
    \item \textbf{距离度量}：使用黎曼距离 $d_R(x, y)$ 替代欧氏距离，能够正确处理流形的几何结构
    \item \textbf{Riemannian centroid}：Robust 估计量基于黎曼中位数而非均值，对异常值更稳健
    \item \textbf{加权 geodesic}：根据数据点到中心的距离分配权重，减少异常值的影响
\end{itemize}

在欧氏空间中，Huber 损失是稳健估计的标准工具；而在黎曼流形上，需要发展相应的几何版本的稳健统计理论。

\subsection{流形上 MCMC 的核心挑战是什么}\label{sec:qa-manifold-mcmc-challenge}

\textbf{问}：在黎曼流形上进行 MCMC 采样的核心挑战是什么？

\textbf{答}：Tang (2024) 关于 Metropolis Adjusted Langevin Algorithm (RMALA) 的工作指出了三大挑战：

\begin{enumerate}
    \item \textbf{曲率适应}：Proposal 分布需要适应局部曲率，否则接受率会很低
    \item \textbf{测地线计算}：Proposal 需要沿着真实测地线而非欧氏直线移动
    \item \textbf{切空间投影}：需要在切空间中进行线性运算，然后投影回流形
\end{enumerate}

RMALA 的核心创新是使用 Riemann Manifold Langevin Dynamics，使得 proposal 分布的协方差与 Fisher 信息矩阵自适应：
$$u' = u - \frac{\epsilon^2}{2} G(x)^{-1} \nabla \log \pi(x) + \epsilon G(x)^{-1/2} \xi$$

其中 $G(x)$ 是 Fisher 信息矩阵（黎曼度量），$\xi \sim N(0, I)$。

\subsection{最小最大率与曲率的关系}\label{sec:qa-minimax-curvature}

\textbf{问}：统计估计的最小最大率（minimax rate）与流形曲率有什么关系？

\textbf{答}：Tang (2022, Annals) 的核心定理建立了这个关系：

对于一个 $d$ 维光滑子流形 $M$，参数估计的 minimax rate 为：
$$\inf_{\hat{\theta}_n} \sup_{\theta \in M} \mathbb{E}_\theta [d_R^2(\hat{\theta}_n, \theta)] \asymp n^{-2/d} + \kappa^2 n^{-4/d}$$

其中 $\kappa$ 是子流形的曲率上界。

关键发现：
\begin{itemize}
    \item $n^{-2/d}$ 项是经典的内点估计率
    \item $\kappa^2 n^{-4/d}$ 项是曲率带来的额外代价
    \item 当 $d$ 较大或曲率较大时，曲率项主导，此时"用简单模型逼近复杂分布"的代价显著
\end{itemize}

几何解释：高曲率区域中，参数空间的微小变化导致分布空间的较大变化，从而增加了估计难度。

\subsection{什么是信息几何视角下的 VAE}\label{sec:qa-info-geometry-vae}

\textbf{问}：从信息几何角度看，VAE 的训练目标是什么？

\textbf{答}：Empirical Bayes VAE (Tang 2021) 从信息几何角度重新审视了 VAE 框架：

VAE 的 ELBO 目标可以解释为在流形上的几何优化：
$$\mathcal{L}(\theta, \phi) = \mathbb{E}_{q_\phi(z|x)}[\log p_\theta(x|z)] - D_{KL}[q_\phi(z|x) \| p(z)]$$

从信息几何看：
\begin{itemize}
    \item $p_\theta(x|z)$ 形成一个指数族分布族构成的流形
    \item $q_\phi(z|x)$ 是近似后验，在另一个流形上
    \item 优化目标是找到这两个流形的最近点
\end{itemize}

潜在问题：Posterior collapse 发生在近似后验流形与先验流形过于接近时，导致解码器忽略潜变量。

\subsection{流形上的两样本检验如何设计}\label{sec:qa-two-sample-test-manifold}

\textbf{问}：如何设计流形上的两样本检验？其功效（power）如何分析？

\textbf{答}：Minimax Two-Sample Test (Tang 2023) 解决了这个问题。设有两个分布 $P$ 和 $Q$ 定义在流形 $\mathcal{M}$ 上，检验问题是：
$$H_0: P = Q \quad \text{vs} \quad H_1: P \neq Q$$

设计思路：
\begin{enumerate}
    \item \textbf{基于距离的检验统计量}：使用黎曼流形上的合适距离（如 Wasserstein 距离的几何版本）
    \item \textbf{基于 Geodesic 的检验}：比较两条样本路径在流形上的积分不变量
    \item \textbf{Spectral 方法}：利用 Laplace-Beltrami 算子的特征函数构造不变量
\end{enumerate}

Minimax 意义下的核心结论：检验的功效下界由分布之间的距离和流形的几何复杂度共同决定。在某些情况下，流形的拓扑不变量（如 Betti 数）会影响检验的可辨别性。

\subsection{什么是 Riemann Manifold MCMC 的接受率理论}\label{sec:qa-rmcmc-acceptance}

\textbf{问}：Riemann Manifold MCMC 的接受率理论是什么？与欧氏空间有何本质区别？

\textbf{答}：RMALA 的接受率公式比欧氏空间版本更复杂，需要考虑流形的几何结构：

在流形上，Metropolis-Hastings 的接受概率为：
$$a(x'|x) = \min\left\{1, \frac{\pi(x') q(x|x') r(x'|x)}{ \pi(x) q(x'|x) r(x|x')}\right\}$$

其中 $r(x'|x)$ 是从 $x$ 到 $x'$ 的测地线 proposal 密度，包含了黎曼度量的 Jacobian 修正。

关键差异：
\begin{itemize}
    \item \textbf{Jacobian 项}：测地线运动的长度依赖于起点和终点
    \item \textbf{度量依赖}：Proposal 的尺度由 Fisher 信息矩阵 $G(x)$ 决定
    \item \textbf{曲率修正}：在高曲率区域需要额外的修正项
\end{itemize}

Tang (2024) 的理论贡献是证明了：当步长 $\epsilon \to 0$ 时，接受率趋向于 $\min\{1, \exp(-\text{KL divergence})\}$，与欧氏情形一致。

\subsection{为什么信息几何是统一框架}\label{sec:qa-info-geometry-unified}

\textbf{问}：为什么说信息几何是理解这系列论文的统一框架？

\textbf{答}：Tang \& Rong 的系列工作可以从信息几何的角度统一理解：

\begin{enumerate}
    \item \textbf{Fisher 信息 = Riemannian 度量}：概率分布空间的自然度量就是 Fisher 信息
    \item \textbf{指数族 = 平坦子流形}：指数族分布构成对偶平坦的流形
    \item \textbf{KL 散度 = 散度几何}：KL 散度是对偶平坦空间中的"距离"
    \item \textbf{推断 = 投影}：贝叶斯后验是数据点向先验流形的投影
    \item \textbf{曲率 = 估计复杂度}：流形的曲率直接影响推断的效率下界
\end{enumerate}

这个框架的核心价值在于：它将概率论、几何、优化和统计推断统一在同一个语言下，使得我们可以系统地理解不同方法之间的关系，并指导新方法的设计。

\subsection{扩散模型的统计推断意义}\label{sec:qa-diffusion-statistical-inference}

\textbf{问}：扩散模型（Score-based models）除了生成任务外，在统计推断中有什么应用？

\textbf{答}：Tang (2024) 揭示了扩散模型在统计推断中的深层应用：

\begin{enumerate}
    \item \textbf{密度估计}：通过学习 score 函数，可以估计数据流形上的概率密度
    \item \textbf{分布检验}：Score 的结构可以用于检测数据是否来自特定流形
    \item \textbf{生成模型的置信区间}：基于扩散模型的采样可以构造参数估计的置信集
    \item \textbf{缺失数据推断}：扩散模型可以补全高维数据中的缺失部分
\end{enumerate}

更深层的联系：Score 函数 $\nabla_x \log p(x)$ 本质上是 Fisher score 的推广形式，因此扩散模型学习的正是概率分布的几何结构。

\subsection{Riemann-Huber 损失函数}\label{sec:qa-riemann-huber-loss}

\textbf{问}：Riemann-Huber 损失函数是什么？它如何推广了欧氏空间的 Huber 损失？

\textbf{答}：Riemann-Huber 损失是 Tang (2023) 在流形上定义的稳健损失函数，形式为：
$$\rho_R(x, y) = \begin{cases} \frac{1}{2} d_R^2(x, y) & \text{if } d_R(x, y) \leq \delta \\ \delta \cdot d_R(x, y) - \frac{\delta^2}{2} & \text{if } d_R(x, y) > \delta \end{cases}$$

其中 $d_R(\cdot, \cdot)$ 是黎曼距离，$\delta$ 是调优参数。

\textbf{与欧氏 Huber 损失的关系}：当流形退化为欧氏空间时，$d_R$ 变为欧氏距离，Riemann-Huber 损失退化为经典 Huber 损失。几何版本的创新在于用黎曼距离替代欧氏距离，使得：
\begin{itemize}
    \item 在正常区域（$d_R \leq \delta$）使用二次损失，保持有效性
    \item 在异常区域（$d_R > \delta$）使用线性损失，控制异常值影响
    \item 整个损失函数在流形上内蕴定义，不依赖坐标选择
\end{itemize}

基于此损失的估计量兼具崩溃点（breakdown point）和有效性。

\subsection{流形结构如何被识别}\label{sec:qa-manifold-structure-identification}

\textbf{问}：Tang (2022) 的工作如何从数据中估计流形的维数和曲率？识别流形结构的实际方法是什么？

\textbf{答}：识别流形结构的核心方法是通过构造 $k$-NN 图（$k$-近邻图）并分析其几何性质：

\begin{enumerate}
    \item \textbf{维数估计}：分析 $k$-NN 距离的缩放行为。设第 $k$ 个近邻距离为 $r_k$，则 $r_k \asymp k^{1/d}$，通过对不同 $k$ 拟合这一关系可以估计 $d$
    \item \textbf{曲率估计}：通过局部 PCA 或局部切空间对齐，估计切空间的弯曲程度，进而得到曲率上界 $\kappa$ 的估计
    \item \textbf{拓扑估计}：利用图的连通性和 Betti 数等拓扑不变量刻画全局结构
\end{enumerate}

几何直觉：$k$-NN 图在流形上是"粗粒度"的离散近似，其几何统计量反映了底层流形的维数和曲率。

\subsection{Fisher 信息与黎曼度量的等价性}\label{sec:qa-fisher-riemann-metric}

\textbf{问}：为什么 Fisher 信息矩阵可以视为概率分布空间上的黎曼度量？

\textbf{答}：Fisher 信息矩阵定义为：
$$g_{ij}(\theta) = \mathbb{E}\left[\frac{\partial \log p(x|\theta)}{\partial \theta_i} \cdot \frac{\partial \log p(x|\theta)}{\partial \theta_j}\right]$$

它满足黎曼度量的两条关键性质：

\begin{itemize}
    \item \textbf{正定性}：对任意非零切向量 $v$，有 $v^T g(\theta) v > 0$
    \item \textbf{对称性}：$g_{ij} = g_{ji}$
\end{itemize}

更深刻的是，Fisher 信息度量具有信息几何的"自然"意义：在参数扰动下，两个分布之间的 KL 散度的二阶近似由 Fisher 信息给出。这使得 Fisher 信息成为概率分布空间上唯一的"自然"黎曼度量——它不依赖于参数化方式。

\subsection{统计推断为何是几何投影}\label{sec:qa-statistical-inference-projection}

\textbf{问}：从信息几何的角度，为什么说各种统计推断问题本质上都是几何投影？

\textbf{答}：信息几何将概率分布空间视为黎曼流形，各种推断任务对应不同类型的投影：

\begin{enumerate}
    \item \textbf{最大似然估计}：将经验分布 $\hat{P}_n$ 投影到模型流形 $\mathcal{M}$ 上最近点。投影方向由 score 函数 $\nabla_\theta \log p(x|\theta)$ 决定
    \item \textbf{贝叶斯推断}：在后验流形与先验流形之间寻找几何平均——后验分布是数据点沿 score 方向的"几何移动"
    \item \textbf{变分推断}：将真实后验流形投影到近似后验流形族上，目标是极小化 KL 散度（流形上的散度几何）
\end{enumerate}

几何直觉：投影是流形上最短路径（测地线）的对偶概念。信息几何提供了统一语言，使得看似不同的推断方法都可以理解为"在流形上找最近点"。

\subsection{RMALA 的 Proposal 动态}\label{sec:qa-rmala-proposal-dynamics}

\textbf{问}：RMALA 的 proposal 分布 $x' = x - \frac{\epsilon^2}{2} G(x)^{-1} \nabla \log \pi(x) + \epsilon G(x)^{-1/2} \xi$ 中，各项的几何意义是什么？

\textbf{答}：这一 dynamics 是 Riemann Manifold Langevin Dynamics，其各项有明确的几何含义：

\begin{itemize}
    \item \textbf{$\nabla \log \pi(x)$（梯度项）}：指向后验密度上升最快的方向——在黎曼流形上，这是切空间中的自然梯度
    \item \textbf{$G(x)^{-1}$（度量逆）}：将自然梯度转换到普通欧氏梯度坐标。由于 Fisher 信息矩阵 $G(x)$ 定义了黎曼度量，需要用其逆来定义恰当的梯度
    \item \textbf{$G(x)^{-1/2} \xi$（随机项）}：在黎曼度量意义下的各向同性噪声。它保证 proposal 分布在黎曼体积形式下是均匀的
    \item \textbf{$\epsilon$（步长）}：控制 proposal 的尺度。当 $\epsilon \to 0$ 时，接受率趋向 $\min\{1, \exp(-\text{KL})\}$
\end{itemize}

与欧氏 LMC 的区别在于：欧氏版本使用固定的单位协方差噪声，而 RMALA 的协方差 $G(x)^{-1}$ 自适应于当前点的局部几何。

\subsection{与 Amari 信息几何框架的关系}\label{sec:qa-amari-framework}

\textbf{问}：Tang \& Rong 的工作与 Amari 的经典信息几何框架有何联系和区别？

\textbf{答}：Tang \& Rong 的工作在 Amari 框架的基础上有重要扩展：

\textbf{继承的核心理论}：
\begin{itemize}
    \item Fisher 信息 = Riemannian 度量（Amari 的基本发现）
    \item 对偶联络（指数联络与混合联络）
    \item 勾股定理在平坦流形上的应用
    \item $\alpha$-联络族
\end{itemize}

\textbf{关键发展}：
\begin{enumerate}
    \item \textbf{从平坦流形扩展到一般弯曲流形}：Amari 的经典理论主要处理对偶平坦的指数族，而 Tang 将理论推广到任意光滑流形，建立了曲率与统计效率的精确关系
    \item \textbf{曲率-效率定理}：这是 Amari 框架中缺失的核心连接——曲率如何影响估计的 minimax 率
    \item \textbf{现代机器学习应用}：将几何方法应用于 VAE、扩散模型等深度生成模型
\end{enumerate}

\textbf{关系}：Amari 框架提供了理论语言和基本工具，Tang \& Rong 则用这个语言建立了更精确的统计效率理论。

\subsection{开放问题与未来研究方向}\label{sec:qa-open-problems}

\textbf{问}：Tang \& Rong 的系列工作之后，流形推断领域有哪些重要开放问题值得关注？

\textbf{答}：基于 Tang (2022-2024) 的开创性贡献，以下方向值得进一步探索：

\begin{enumerate}
    \item \textbf{高维流形推断}：当流形维数 $d$ 随样本量 $n$ 增长时（如 $d \asymp n^\alpha$），minimax 率如何刻画？这需要全新的分析技术
    \item \textbf{非光滑流形}：现有理论假设流形光滑（$C^\infty$），但实际数据可能存在边缘、尖点等奇异几何结构。如何建立非光滑流形上的统计理论？
    \item \textbf{计算-统计权衡}：现有理论给出了最优率，但对应的最优算法是否可计算？如何在保持理论性质的同时设计有效算法？
    \item \textbf{神经网络表示空间}：深度学习模型（如 transformers、diffusion models）学习到的表示空间是否构成有意义的黎曼流形？如何用几何工具分析这些空间？
    \item \textbf{因果推断与几何}：因果结构的识别是否可以借助流形方法？因果流形与统计流形有何本质区别？
\end{enumerate}

其中因果推断与几何的结合是最前沿的方向之一，可能带来因果推断的新突破。
